\begin{titlepage}

	\parbox{2cm}{
	\begin{picture}(18,4)
	    \put(-21,230){\includegraphics[width=2cm,height=2cm]{imagen/IPN.png}}
	    \put(0,-200){\includegraphics[width=0.5cm,height=15.3cm]{imagen/LineaAzul.png}}
	    \put(-37,-275){\includegraphics[width=3cm,height=2.5cm]{imagen/logoescom.png}}
	\end{picture}}
	\parbox{14cm}{\vspace{1cm} 
	\begin{center}
	    {\fontsize{20}{30}\textbf{ INSTITUTO POLITÉCNICO NACIONAL}}\\[0.2cm]
	    
	    {\fontsize{16}{20} \textbf{ESCUELA SUPERIOR DE CÓMPUTO}}\vspace{1cm}\\[1cm]
	    {\fontsize{20}{20} \textbf{ESCOM}}\vspace{2cm}\\
	    
	    {\fontsize{14}{20} \textit{Trabajo Terminal}}\vspace{1cm}\\
	    {\fontsize{16}{20} \textbf{``Yolotl: un videojuego para fomentar la cultura''}}\vspace{0.5cm}\\
	    {\fontsize{14}{20} \textit{2017-A035}}\vspace{1.5cm}\\
	    {\fontsize{14}{20} \textit{Presentan}}\\
	    {\fontsize{14}{20} \textbf{Hernández Bautista Yasmine Pilar}}\vspace{1cm}\\
	   	{\fontsize{14}{20} \textbf{Márquez Hernández Karla Rocío}}\vspace{1.5cm}\\
	   \fontsize{14}{20} \textit{Directores}\\

	    
	    
{\fboxrule=0pt \fboxsep=6pt	    
\fbox{{\fontsize{14}{20} \textbf{\textit{M. en C. Rafael Norman Saucedo Delgado}}.}\vspace{3.5cm}}
\fbox{{\fontsize{14}{20} \textbf{\textit{Lic. Ulises Vélez Saldaña}}.}\vspace{3.5cm}}}\\[3.5cm]
\end{center}
	    
	    \hfill  \fontsize{12}{20} \textit{Junio 2018}
	    
	}
\end{titlepage}

\thispagestyle{empty}

\parbox{18cm}{
\parbox{1.5cm}{
\includegraphics[width=1.5cm,height=2.5cm]{imagen/IPN.png}
}
\parbox{12cm}{
\centering{{\fontsize{16}{0}\textbf{ INSTITUTO POLITÉCNICO NACIONAL}}\\	
{\fontsize{16}{0} \textbf{ESCUELA SUPERIOR DE CÓMPUTO}}}
}
\parbox{1.5cm}{
\includegraphics[width=2.5cm,height=2cm]{imagen/logoescom.png}
}\vspace{1.5cm} 
}

	\begin{center}
	    
	    
\begin{multicols}{2} 
\raggedright{{\fontsize{14}{20} No. de TT:2017-A035}} 

\raggedleft{{\fontsize{14}{20} 31 de junio de 2018}}
\end{multicols}\vspace{1cm}

	    {\fontsize{14}{20} \textit{Documento Técnico Parte A}}\vspace{1cm}\\
	    {\fontsize{16}{20} \textbf{``Yolotl: un videojuego para fomentar la cultura''}}\vspace{1.5cm}\\
	    {\fontsize{14}{20} \textit{Presentan}}\\
	    {\fontsize{14}{20} \textbf{Hernández Baustista Yasmine Pilar\footnote{daughterofthewind10@gmail.com}}}\vspace{1cm}\\
	    {\fontsize{14}{20} \textbf{Márquez Hernández Karla Rocío\footnote{winged.zelda@gmail.com}}}\vspace{1cm}\\
	   
	   \fontsize{14}{20} \textit{Directores}\vspace{1.5cm}\\
	    
	    
{\fboxrule=0pt \fboxsep=12pt	    
\fbox{{\fontsize{14}{20} \textbf{\textit{M. en C. Rafael Norman Saucedo Delgado}}.}\vspace{5.5cm}}
\fbox{{\fontsize{14}{20} \textbf{\textit{Lic. Ulises V\'elez Saldaña}}.}
}}
\end{center}
\begin{center}
{\fontsize{14}{20} \textbf{RESUMEN}}
\end{center}
En México la industria de videojuegos tiene una alta demanda de consumo; sin embargo, existen pocos estudios que desarrollen videojuegos basados en la cultura mexicana. Actualmente en México existe un fuerte desinterés en la cultura nacional. El presente trabajo terminal consiste en el desarrollo de un videojuego que fomente la cultura con temática de la cultura mexica. 
\\
 
\textbf{\textit{Palabras clave}}. –  Cultura mexica, desarrollo tecnológico, ingeniería de software, videojuego.
